\documentclass[11pt]{article}
\usepackage[T1]{fontenc}
\usepackage{lmodern}
\usepackage[margin=1in]{geometry}
\usepackage{amsmath,amssymb,amsthm,mathtools}
\usepackage{microtype}
\usepackage{enumitem}
\usepackage{tikz}
\usepackage[colorlinks=true,linkcolor=blue!45!black,citecolor=green!35!black,urlcolor=blue!55!black]{hyperref}
\usepackage[all]{hypcap}
\hypersetup{
  pdftitle={A product counterexample to the simplex maximum for projection-body volume},
  pdfauthor={Rewritten for readability from the OpenAI manuscript of the same title},
  pdfsubject={An elementary counterexample in dimension twenty}
}
\newtheorem{theorem}{Theorem}
\newtheorem{lemma}[theorem]{Lemma}
\newtheorem{proposition}[theorem]{Proposition}
\newtheorem{corollary}[theorem]{Corollary}
\newcommand{\R}{\mathbb R}
\DeclareMathOperator{\conv}{conv}
\DeclareMathOperator{\vol}{vol}
\DeclareMathOperator{\proj}{proj}
\DeclareMathOperator{\mult}{mult}
\newcommand{\unverified}[1]{\par\smallskip\noindent\fbox{\parbox{0.95\linewidth}{\textbf{Unverified step.} #1}}\par\smallskip}
\newcommand{\step}[1]{\par\smallskip\noindent\textbf{#1}\ }
\title{A product counterexample to the simplex maximum\\
for projection-body volume}
\author{Rewritten for readability from the OpenAI manuscript of the same title}
\date{September 24, 2026}
\raggedbottom

\begin{document}
\maketitle

\begin{abstract}
The product of two ten-dimensional simplices has larger normalized
projection-body volume than a twenty-dimensional simplex. This gives
a counterexample to Brannen's proposed simplex maximum.
\end{abstract}

\section{The result and how it is proved}

\subsection{Setting and definitions}

Throughout, \(d\ge2\). A \emph{convex body} in \(\R^d\) is a compact convex
set with nonempty interior. We write \(|K|\) for its \(d\)-dimensional
volume, \(\langle\cdot,\cdot\rangle\) and \(\|\cdot\|\) for the Euclidean
inner product and norm, \(S^{d-1}\) for the unit sphere, \(B_2^j\) for the
Euclidean unit ball of \(\R^j\), and \(\kappa_j=|B_2^j|\). For \(u\in S^{d-1}\),
\(\proj_{u^\perp}\) is the orthogonal projection onto the hyperplane
\(u^\perp\). The \emph{support function} of a compact convex set \(C\) is
\(h_C(x)=\max_{y\in C}\langle x,y\rangle\). The \emph{Minkowski sum} is
\(C+D=\{c+d:c\in C,\ d\in D\}\), and \([p,q]\) denotes the segment from \(p\)
to \(q\).

\smallskip\noindent\textbf{Projection body.}
For a convex body \(K\subset\R^d\), the projection body \(\Pi K\) (also
\(\Pi_dK\)) is the compact convex set whose support function is
\[
 h_{\Pi K}(u)=\vol_{d-1}\bigl(\proj_{u^\perp}K\bigr)\quad(u\in S^{d-1}),
 \qquad h_{\Pi K}(tu)=t\,h_{\Pi K}(u)\quad(t\ge0).
\]
This is the definition throughout. A compact convex set is determined by
its support function (fact (S4) in Section~\ref{sec:facets}), so \(\Pi K\)
is unique if it exists. Existence is proved for polytopes in
Lemma~\ref{lem:facets}, with an explicit formula, and for arbitrary convex
bodies in Step~6 of the proof of Lemma~\ref{lem:affine}. The main theorem
and the corollary involve polytopes only.

\smallskip\noindent\textbf{The functional and the simplex constant.}
\begin{equation}\label{eq:functional}
 R_d(K)=\frac{|\Pi K|}{|K|^{d-1}},
 \qquad c_d=\frac{(d+1)d^d}{d!}.
\end{equation}
Proposition~\ref{prop:simplex} shows that \(R_d(T)=c_d\) for every
\(d\)-simplex \(T\). Brannen conjectured in 1996 that simplices maximize
\(R_d\) among all convex bodies~\cite{Brannen}. In view of the simplex
value, this is the proposed bound
\begin{equation}\label{eq:proposed-bound}
 |\Pi K|\le c_d|K|^{d-1}\qquad\text{for every convex body }K\subset\R^d.
\end{equation}

\subsection{Main results}

\begin{theorem}\label{thm:counterexample}
Let \(T_{10}=\conv(0,e_1,\ldots,e_{10})\subset\R^{10}\), where the
\(e_i\) are the coordinate unit vectors, and let
\(K=T_{10}\times T_{10}\subset\R^{20}\). Then
\[
 \frac{R_{20}(K)}{c_{20}}
 =\frac{121\binom{20}{10}}{21\cdot2^{20}}
 =\frac{22\,355\,476}{22\,020\,096}>1.
\]
Thus \(K\) violates \eqref{eq:proposed-bound} in dimension twenty.
\end{theorem}

The second result, Corollary~\ref{cor:exponential-excess} in
Section~\ref{sec:counterexample}, says that the excess over the simplex
value can be made exponential in the dimension: there are \(\lambda>1\) and
\(N\) such that for every \(n\ge N\) some product of simplices
\(K_n\subset\R^n\) has \(R_n(K_n)\ge\lambda^n c_n\).

\subsection{Context (reported from the source; not used in any proof)}

Everything in this subsection is taken from the source's introduction. It
rests on the cited papers, which this document does not check, and no
proof below uses it.

Petty established the affine covariance of the projection-body
operator~\cite{Petty}; see also \cite[Equation~(1)]{Ludwig}. The functional
\(R_d\) involves the volume of \(\Pi K\) itself, not of its polar; the
classical projection inequalities involving the polar concern a different
functional. For \(d\ge3\), Petty's separate minimum problem asks whether
ellipsoids are the only minimizers of the same functional \(R_d\); see
\cite[Section~1]{Saroglou} for these distinctions.

In dimension three, Chen, Feng, Li, Xi, and Xu settle the minimum problem
and also prove \(R_3(K)\le18=c_3\) for all three-dimensional convex bodies,
with tetrahedra attaining equality~\cite[Theorems~1.1 and~1.2]{ChenEtAl}.
Feng, Hu, Liu, and Xu give counterexamples to \eqref{eq:proposed-bound} in
every dimension \(d\ge9\), using polytopes with at most \(d+2\)
facets~\cite[Theorem~1.3]{FengEtAl}; their construction uses perturbed
projections of a cube and Minkowski's existence
theorem~\cite[Section~4]{FengEtAl}. Taken at face value, that cited theorem
already shows that \eqref{eq:proposed-bound} fails in dimension twenty. So
the failure of \eqref{eq:proposed-bound} is not new with the present paper.
The source says so: ``Their result already gives a negative answer to the
unrestricted-dimensional simplex-maximum question.'' It describes its own
contribution as ``a direct Cartesian-product construction and its exact
value: two ten-dimensional simplices suffice''.
Theorem~\ref{thm:counterexample} is proved below without using any cited
result. The optimal upper constant and its maximizing bodies remain
separate questions.

For centrally symmetric bodies, Henk gives the lower bound
\(2^d(9/8)^{\lfloor d/3\rfloor}\) for the optimal upper
constant~\cite{Henk}. This sentence is the source's. In its context,
``upper constant'' refers to upper bounds for \(R_d\), here over centrally
symmetric bodies. The source quotes it only as background on how such
constants grow with \(d\), and it draws no comparison with
Corollary~\ref{cor:exponential-excess}. That corollary measures growth
\emph{relative to the simplex value} \(c_n\) only.

\subsection{Proof at a glance}

\textbf{Key idea.} For polytopes, \(R\) is multiplicative under Cartesian
products:
\[
 R_{r+s}(A\times B)=R_r(A)\,R_s(B)\qquad(r,s\ge2).
\]
The simplex constants are not: \(c_{10}^2>c_{20}\). Hence
\(R_{20}(T_{10}\times T_{10})=c_{10}^2>c_{20}\). In general, call \(C\) a
\emph{universal upper constant in dimension \(d\)} if \(R_d(K)\le C\) for
all convex bodies \(K\subset\R^d\). Since \(T_r\times T_s\) is a convex
body in \(\R^{r+s}\) with \(R_{r+s}=c_rc_s\), every such constant in
dimension \(r+s\) is at least \(c_rc_s\). The conjecture would need
\(c_rc_s\le c_{r+s}\).

An informal heuristic, not used in any proof, explains why this fails.
Stirling's formula \(d!\approx\sqrt{2\pi d}\,(d/e)^d\) gives
\(c_d\approx(d+1)e^d/\sqrt{2\pi d}\), hence
\[
 \frac{c_rc_s}{c_{r+s}}\approx
 \frac{(r+1)(s+1)}{r+s+1}\sqrt{\frac{r+s}{2\pi rs}}.
\]
The exponential factors have cancelled. For \(r=s\) the right side is
\((r+1)^2/((2r+1)\sqrt{\pi r})\), which grows like \(\sqrt r/(2\sqrt\pi)\).
At \(r=10\) it is about \(1.03\); the exact value, from
Theorem~\ref{thm:counterexample}, is about \(1.015\).

\textbf{Steps.}
\begin{enumerate}[leftmargin=2em,itemsep=2pt]
\item \emph{Facet formula (Lemma~\ref{lem:facets}).} For a polytope \(P\),
 \(\Pi P\) is a sum of segments, one per facet, each spanned by the facet's
 unit normal scaled by its area. Everything else is computed from this.
\item \emph{Product identity (Proposition~\ref{prop:product}).} The facets
 of \(A\times B\) are \(F\times B\) and \(A\times G\). Step~1 then gives
 \(\Pi(A\times B)=(|B|\Pi A)\times(|A|\Pi B)\), and taking volumes gives
 multiplicativity.
\item \emph{Affine covariance (Lemma~\ref{lem:affine}).}
 \(R_d(LP)=R_d(P)\) for invertible linear \(L\). It is used only to pass
 from the standard simplex to all simplices.
\item \emph{Simplex value (Proposition~\ref{prop:simplex}).} By Step~1,
 \(\Pi T_d\) is a scaled translate of ``unit cube plus one segment'', whose
 volume is \(d+1\). This gives \(R_d(T_d)=c_d\).
\item \emph{Arithmetic.} Theorem~\ref{thm:counterexample} is the exact
 evaluation of \(c_{10}^2/c_{20}\).
 Corollary~\ref{cor:exponential-excess} iterates the product identity and
 uses the bound \(c_n<(n+1)(11/4)^n\).
\end{enumerate}

\begin{center}
\begin{tikzpicture}[every node/.style={draw,rounded corners,font=\small,align=center,inner sep=4pt},>=stealth]
 \node (L2) at (0,0) {Lemma~\ref{lem:facets}\\facet formula};
 \node (P3) at (4.2,1.1) {Prop.~\ref{prop:product}\\product identity};
 \node (L4) at (4.2,-1.1) {Lemma~\ref{lem:affine}\\affine covariance};
 \node (P5) at (8.6,-1.1) {Prop.~\ref{prop:simplex}\\simplex value};
 \node (T1) at (12.8,1.1) {Theorem~\ref{thm:counterexample}};
 \node (C6) at (12.8,-1.1) {Corollary~\ref{cor:exponential-excess}};
 \draw[->] (L2) -- (P3);
 \draw[->] (L2) -- (L4);
 \draw[->] (L4) -- (P5);
 \draw[->] (L2.south) -- ++(0,-1.9) -| (P5.south);
 \draw[->] (P3) -- (T1);
 \draw[->] (P5) -- (T1);
 \draw[->] (P3) -- (C6);
 \draw[->] (P5) -- (C6);
\end{tikzpicture}
\end{center}

\noindent Two remarks for a reader short of time.
\begin{itemize}[leftmargin=2em,itemsep=2pt]
\item Theorem~\ref{thm:counterexample} and
 Corollary~\ref{cor:exponential-excess} use only the \emph{standard}
 simplex part of Proposition~\ref{prop:simplex}, which does not need
 Lemma~\ref{lem:affine}. Lemma~\ref{lem:affine}, in its polytope case, is
 needed only for the sentence ``every simplex has value \(c_d\)'', which
 identifies \(c_d\) as the simplex value in Brannen's conjecture.
\item The second half of Lemma~\ref{lem:affine} (arbitrary convex bodies,
 Steps~5--8) is used nowhere else. It is kept because the source includes
 it and because it shows that \(\Pi K\) exists for every convex body.
\end{itemize}

\subsection{Notation changes relative to the source}

Four symbols are renamed to avoid clashes, with no change in the
mathematics. In the proof of Lemma~\ref{lem:affine}, the outer radius \(R\)
becomes \(\rho\) (the letter \(R_d\) is the functional), the factors
\(\lambda_m\) become \(\mu_m\) (\(\lambda\) is the constant of
Corollary~\ref{cor:exponential-excess}), and the functions \(F,f_m\) become
\(\varphi,\varphi_m\) (\(F\) is a facet). In the proof of
Corollary~\ref{cor:exponential-excess}, the constant \(A\) becomes
\(\alpha\) (\(A\) is a polytope in Proposition~\ref{prop:product}).

\section{Facet vectors and Cartesian products}\label{sec:facets}

\subsection{Four facts about support functions}

Let \(C,D\) be compact convex sets. Each of (S1)--(S3) is immediate from
\(h_C(x)=\max_{y\in C}\langle x,y\rangle\).
\begin{itemize}[leftmargin=3em,itemsep=2pt]
\item[(S1)] \(h_{C+D}=h_C+h_D\), and \(h_{tC}=t\,h_C\) for \(t\ge0\).
 The maximum of \(\langle x,c+d\rangle\) over independent \(c\in C\),
 \(d\in D\) splits into two maxima.
\item[(S2)] For \(C\subset\R^r\), \(D\subset\R^s\) and
 \((u,v)\in\R^r\times\R^s\): \(h_{C\times D}(u,v)=h_C(u)+h_D(v)\), for the
 same reason. For a linear map \(M\): \(h_{MC}(u)=h_C(M^tu)\), because
 \(\langle u,Mc\rangle=\langle M^tu,c\rangle\).
\item[(S3)] \(h_{[-v/2,\,v/2]}(u)=\tfrac12|\langle v,u\rangle|\), since a
 linear function on a segment is maximal at an endpoint. Also
 \(h_{tB_2^d}(u)=t\|u\|\).
\item[(S4)] If \(h_C\le h_D\) on \(\R^d\), then \(C\subset D\). In
 particular \(h_C=h_D\) implies \(C=D\).
\end{itemize}
\emph{Proof of (S4).} Suppose \(x\in C\setminus D\). Let \(y\in D\) be the
point of \(D\) nearest to \(x\), so \(x\ne y\). For \(z\in D\) and
\(t\in[0,1]\) the point \(y+t(z-y)\) lies in \(D\), so
\(g(t)=\|x-y-t(z-y)\|^2\) is minimal at \(t=0\). Hence
\(g'(0)=-2\langle x-y,z-y\rangle\ge0\). Thus
\(\langle x-y,z\rangle\le\langle x-y,y\rangle\) for all \(z\in D\), that
is, \(h_D(x-y)=\langle x-y,y\rangle\). But
\(\langle x-y,x\rangle=\langle x-y,y\rangle+\|x-y\|^2>h_D(x-y)\), while
\(x\in C\) gives \(\langle x-y,x\rangle\le h_C(x-y)\le h_D(x-y)\). This is
a contradiction. \qed

\smallskip
Support functions are positively homogeneous (\(h_C(tx)=t\,h_C(x)\) for
\(t\ge0\)). So an identity or inequality between support functions need
only be checked on \(S^{d-1}\).

\subsection{The facet formula}

We use the following standard facts about a full-dimensional convex
polytope \(P\subset\R^d\). It has finitely many faces, and its
\emph{facets} are the faces of dimension \(d-1\). For a facet \(F\), write
\(s_F\) for its \((d-1)\)-dimensional volume, \(\nu_F\) for its outward
unit normal, and \(H_F=\{x:\langle\nu_F,x\rangle=h_P(\nu_F)\}\) for its
hyperplane. Then:
\begin{itemize}[leftmargin=3em,itemsep=2pt]
\item[(P1)] \(P=\bigcap_F\{x:\langle\nu_F,x\rangle\le h_P(\nu_F)\}\) and
 \(F=P\cap H_F\).
\item[(P2)] The boundary of \(P\) is the union of the facets.
\item[(P3)] A point of a facet \(F\) outside the relative interior of
 \(F\), and a point lying in two distinct facets, lies in a face of
 dimension at most \(d-2\).
\end{itemize}
We call \(s_F\nu_F\) the \emph{area-normal vector} of \(F\).

The next lemma is the polytopal case of Cauchy's projection formula; see
\cite[Section~1]{Saroglou}. The proof is included to fix the factor
\(1/2\), on which all later constants depend. The lemma is used in
Proposition~\ref{prop:product}, Lemma~\ref{lem:affine} and
Proposition~\ref{prop:simplex}.

\begin{lemma}\label{lem:facets}
Let \(P\subset\R^d\), \(d\ge2\), be a full-dimensional convex polytope.
Then \(\Pi P\) exists,
\begin{equation}\label{eq:facet}
 h_{\Pi P}(u)=\frac12\sum_Fs_F|\langle\nu_F,u\rangle|,
 \qquad u\in\R^d,
\end{equation}
and therefore
\begin{equation}\label{eq:zonotope}
 \Pi P=\sum_F
 \left[-\frac{s_F\nu_F}{2},\frac{s_F\nu_F}{2}\right],
\end{equation}
where both sums run over the facets \(F\) of \(P\).
\end{lemma}

\begin{proof}
\emph{Plan.} Fix a unit vector \(u\) and let \(\pi=\proj_{u^\perp}\). We
show that \(\pi\) scales areas on \(H_F\) by \(|\langle\nu_F,u\rangle|\)
(Step~1). We then show that almost every point of \(\pi(P)\) is covered by
the projections of exactly two facets (Step~2). Counting the covered area
in two ways gives \eqref{eq:facet} for unit \(u\) (Step~3). The general
case and \eqref{eq:zonotope} follow from (S1), (S3), (S4) (Step~4).

\step{Step 1 (Jacobian).} \emph{For every facet \(F\) and measurable
\(S\subset H_F\),
\(\vol_{d-1}(\pi(S))=|\langle\nu_F,u\rangle|\vol_{d-1}(S)\).}
The restriction of \(\pi\) to \(H_F\) is affine with linear part
\(J\colon\nu_F^\perp\to u^\perp\), \(Jv=v-\langle v,u\rangle u\). So the
claim is \(|\det J|=|\langle\nu_F,u\rangle|\); when this is \(0\), both
sides vanish. Let \(v_1,\dots,v_{d-1}\) be an orthonormal basis of
\(\nu_F^\perp\). Then \(|\det J|\) is the \((d-1)\)-volume of the
parallelepiped spanned by \(Jv_1,\dots,Jv_{d-1}\) in \(u^\perp\). Since
\(u\) is a unit vector orthogonal to all \(Jv_i\), this equals the
\(d\)-volume of the parallelepiped with the extra edge \(u\):
\[
 |\det J|=|\det(Jv_1,\dots,Jv_{d-1},u)|
 =|\det(v_1,\dots,v_{d-1},u)|
 =|\langle u,\nu_F\rangle|\,|\det(v_1,\dots,v_{d-1},\nu_F)|
 =|\langle u,\nu_F\rangle|.
\]
The second equality adds multiples of the last column to the others. The
third expands
\(u=\langle u,\nu_F\rangle\nu_F+\sum_i\langle u,v_i\rangle v_i\) in the
last column, where the terms with a repeated column vanish. The last holds
because \((v_1,\dots,v_{d-1},\nu_F)\) is an orthonormal basis.

\step{Step 2 (generic lines meet exactly two facets).} For
\(y\in u^\perp\) let \(\mult(y)\) be the number of facets \(F\) with
\(y\in\pi(F)\). Clearly \(\mult(y)=0\) for \(y\notin\pi(P)\). \emph{We
claim \(\mult(y)=2\) for all \(y\in\pi(P)\setminus E\), where
\(E\subset u^\perp\) has \((d-1)\)-dimensional measure zero.}

\emph{The exceptional set.} Let \(E\) be the union of \(\pi(G)\) over all
faces \(G\) of dimension at most \(d-2\), together with \(\pi(F)\) over all
facets \(F\) with \(\langle\nu_F,u\rangle=0\). Each \(\pi(G)\) lies in the
image of an affine space of dimension at most \(d-2\). If
\(\langle\nu_F,u\rangle=0\), then \(H_F\) contains the direction \(u\), so
\(\pi(H_F)\) has dimension \(d-2\). Thus \(E\) is a finite union of subsets
of affine subspaces of \(u^\perp\) of dimension at most \(d-2\), hence
null.

\emph{The count.} Let \(y\in\pi(P)\setminus E\) and
\(\ell=y+\R u\); every point of \(\ell\) projects to \(y\). By compactness
and convexity, \(\ell\cap P\) is a nonempty segment \([p,q]\), possibly
with \(p=q\). Both endpoints lie on the boundary of \(P\), because \(\ell\)
leaves \(P\) there.
\begin{enumerate}[label=(\roman*),leftmargin=2.5em,itemsep=2pt]
\item By (P2), \(p\) lies in a facet \(F_p\). If \(p\) lay in a face \(G\)
 of dimension at most \(d-2\), then \(y=\pi(p)\in\pi(G)\subset E\), which
 is excluded. So by (P3), \(p\) is in the relative interior of \(F_p\) and
 in no other facet. Also \(\langle\nu_{F_p},u\rangle\ne0\): otherwise
 \(y=\pi(p)\in\pi(F_p)\subset E\). The same holds for \(q\), with a facet
 \(F_q\).
\item By (P1) and (i), every facet inequality other than that of \(F_p\) is
 strict at \(p\). So in a small ball around \(p\), \(P\) coincides with the
 half-space \(\{\langle\nu_{F_p},x\rangle\le h_P(\nu_{F_p})\}\). As
 \(\langle\nu_{F_p},u\rangle\ne0\), the points \(p+tu\) with \(|t|\) small,
 \(t\ne0\) and \(t\langle\nu_{F_p},u\rangle<0\) are interior points of
 \(P\). Fix one of them, \(x_0\). It lies in \(\ell\cap P=[p,q]\) and
 \(x_0\ne p\), so \(p\ne q\). Also \(x_0\ne q\), because \(q\) is a
 boundary point. Thus \(x_0\) lies strictly between \(p\) and \(q\).
\item If \(x_0\) is interior to \(P\) and \(z\in P\), every point of
 \([x_0,z]\) other than \(z\) is interior to \(P\). (If a ball of radius
 \(\varepsilon\) about \(x_0\) lies in \(P\), convexity puts a ball of
 radius \((1-t)\varepsilon\) about \((1-t)x_0+tz\) in \(P\).) Apply this
 with \(z=p\) and \(z=q\). Since \(x_0\) is strictly between \(p\) and
 \(q\), the two half-open segments cover the open segment between \(p\) and
 \(q\), which therefore lies in the interior of \(P\). So \(\ell\) meets
 the boundary of \(P\) only at \(p\) and \(q\).
\item A facet \(F\) with \(y\in\pi(F)\) contains a point of \(\ell\). This
 point is a boundary point of \(P\), hence is \(p\) or \(q\) by (iii), so
 \(F\in\{F_p,F_q\}\) by (i). Conversely \(y\in\pi(F_p)\) and
 \(y\in\pi(F_q)\). Finally \(F_p\ne F_q\). The function
 \(t\mapsto\langle\nu_{F_p},y+tu\rangle\) has slope
 \(\langle\nu_{F_p},u\rangle\ne0\), so \(H_{F_p}\) meets \(\ell\) in exactly
 one point, namely \(p\). If \(F_q=F_p\), then
 \(q\in H_{F_p}\cap\ell=\{p\}\), contradicting \(p\ne q\).
\end{enumerate}
Hence \(\mult(y)=2\).

\emph{Remark.} The source's exceptional set also contains the relative
boundary of \(\pi(P)\). Discarding it is harmless but not needed. By (ii),
for \(y\in\pi(P)\setminus E\) the line \(\ell\) passes through an interior
point of \(P\), so \(y\) is an interior point of \(\pi(P)\) in \(u^\perp\).

\step{Step 3 (double counting).} Integrating the indicator functions of
the sets \(\pi(F)\) over \(u^\perp\) and using Steps~1 and~2,
\[
 \sum_F s_F|\langle\nu_F,u\rangle|
 =\sum_F\vol_{d-1}(\pi(F))
 =\int_{u^\perp}\mult(y)\,dy
 =2\vol_{d-1}(\pi(P)).
\]
So \(\vol_{d-1}(\proj_{u^\perp}P)=\tfrac12\sum_Fs_F|\langle\nu_F,u\rangle|\)
for every unit vector \(u\).

\step{Step 4 (general \(u\) and the zonotope).} Let \(Z\) be the right side
of \eqref{eq:zonotope}, a compact convex set. By (S1) and (S3),
\(h_Z(u)=\tfrac12\sum_Fs_F|\langle\nu_F,u\rangle|\) for all \(u\in\R^d\).
This function is positively homogeneous of degree one, and by Step~3 it
equals \(\vol_{d-1}(\proj_{u^\perp}P)\) on unit vectors. That is exactly
the definition of the support function of \(\Pi P\): the projection onto
\(u^\perp\) enters only for unit \(u\), and elsewhere \(h_{\Pi P}\) is
defined by homogeneity. So \(\Pi P\) exists and equals \(Z\), uniquely by
(S4), and \eqref{eq:facet} holds for all \(u\in\R^d\).
\end{proof}

\subsection{The product identity}

The next proposition is the mechanism behind both main results. It is used
in the proofs of Theorem~\ref{thm:counterexample} and
Corollary~\ref{cor:exponential-excess}. The hypothesis \(r,s\ge2\) is there
only because projection bodies are defined in dimension at least two.

\begin{proposition}\label{prop:product}
Let \(A\subset\R^r\) and \(B\subset\R^s\) be full-dimensional convex
polytopes, where \(r,s\ge2\). In the orthogonal product space
\(\R^{r+s}=\R^r\times\R^s\),
\begin{equation}\label{eq:product-body}
 \Pi(A\times B)=(|B|\Pi A)\times(|A|\Pi B).
\end{equation}
In particular,
\begin{equation}\label{eq:multiplicative}
 R_{r+s}(A\times B)=R_r(A)R_s(B).
\end{equation}
\end{proposition}

\begin{proof}
\emph{Plan.} Find the facets of \(A\times B\) with their normals and areas
(Steps~1--2), insert them into the facet formula (Step~3), and take volumes
(Step~4). As in the source, the bare letter \(s\) is the dimension of
\(B\), while \(s_F,s_G\) (always with a subscript) are facet areas.

\step{Step 1 (facets of \(A\times B\)).} \emph{The facets of \(A\times B\)
are the sets \(F\times B\), with outward unit normal \((\nu_F,0)\), for
facets \(F\) of \(A\), and the sets \(A\times G\), with outward unit normal
\((0,\nu_G)\), for facets \(G\) of \(B\).}

For \(u\in\R^r\) let \(A_u\) be the set of maximizers of
\(\langle u,\cdot\rangle\) on \(A\), and define \(B_v\) likewise. Thus
\(A_0=A\). For \(u\ne0\), \(A_u\) lies in the hyperplane
\(\{x:\langle u,x\rangle=h_A(u)\}\), so \(\dim A_u\le r-1\). Let \(\Phi\)
be a facet of \(A\times B\) with outward unit normal \((u,v)\). By (P1),
\(\Phi\) is the set of maximizers of
\((a,b)\mapsto\langle u,a\rangle+\langle v,b\rangle\) on \(A\times B\). The
two terms are maximized independently, so \(\Phi=A_u\times B_v\). The
affine hull of a product is the product of the affine hulls, so
\(\dim\Phi=\dim A_u+\dim B_v\). If \(u\ne0\ne v\), this is at most
\(r+s-2\), contradicting \(\dim\Phi=r+s-1\). So exactly one of \(u,v\)
vanishes.

Suppose \(v=0\). Then \(B_v=B\) and \(\dim A_u=r-1\), so \(A_u\) is a facet
\(F\) of \(A\), and \(\|u\|=1\). The hyperplane
\(\{x:\langle u,x\rangle=h_A(u)\}\) contains the \((r-1)\)-dimensional set
\(F\), so it equals \(H_F\), and \(u=\pm\nu_F\). Both \(u\) and \(\nu_F\)
have \(A\) on the side where their functional is at most its value on
\(H_F\), and \(A\not\subset H_F\). Hence \(u=\nu_F\). The case \(u=0\) is
symmetric. Conversely, \(F\times B\) is the set of maximizers of
\((\nu_F,0)\) and has dimension \(r+s-1\), so it is a facet; likewise
\(A\times G\).

\step{Step 2 (areas).}
\(\vol_{r+s-1}(F\times B)=s_F|B|\) and \(\vol_{r+s-1}(A\times G)=|A|s_G\).
Indeed, \(F\times B\) lies in \(H_F\times\R^s\). Because the two coordinate
spaces are orthogonal, an isometry \(H_F\cong\R^{r-1}\) times the identity
of \(\R^s\) is an isometry \(H_F\times\R^s\cong\R^{r-1}\times\R^s\), and
Lebesgue measure on the latter is the product measure. Hence the
area-normal vectors of the facets of \(A\times B\) are
\[
 (|B|s_F\nu_F,0)
 \quad\hbox{and}\quad
 (0,|A|s_G\nu_G).
\]

\step{Step 3 (the projection body).} By \eqref{eq:facet} for \(A\times B\),
then for \(A\) and for \(B\),
\[
 h_{\Pi(A\times B)}(u,v)
 =\frac12\sum_F|B|s_F|\langle\nu_F,u\rangle|
  +\frac12\sum_G|A|s_G|\langle\nu_G,v\rangle|
 =|B|h_{\Pi A}(u)+|A|h_{\Pi B}(v).
\]
By (S1) and (S2), the right side is the support function of
\((|B|\Pi A)\times(|A|\Pi B)\) at \((u,v)\). So (S4) gives
\eqref{eq:product-body}.

\step{Step 4 (volumes).} Volume is multiplicative on products, and scaling
by \(t>0\) multiplies \(j\)-dimensional volume by \(t^j\). Since
\(\Pi A\subset\R^r\) and \(\Pi B\subset\R^s\), we get
\(|\Pi(A\times B)|=|B|^r|\Pi A|\cdot|A|^s|\Pi B|\). Also
\(|A\times B|=|A||B|\). Hence
\[
 R_{r+s}(A\times B)
 =\frac{|B|^r|A|^s|\Pi A||\Pi B|}{(|A||B|)^{r+s-1}}
 =\frac{|\Pi A|}{|A|^{r-1}}\cdot\frac{|\Pi B|}{|B|^{s-1}},
\]
because \(|A|^{s}/|A|^{r+s-1}=|A|^{-(r-1)}\) and
\(|B|^{r}/|B|^{r+s-1}=|B|^{-(s-1)}\).
\end{proof}

\subsection{Affine covariance}

To use simplices as a benchmark we need all \(d\)-simplices to have the
same value of \(R_d\). The next lemma supplies this. Its polytope case
(Steps~1--4) is used once, in the last step of
Proposition~\ref{prop:simplex}. Its general case (Steps~5--8) is used
nowhere else.

\begin{lemma}\label{lem:affine}
Let \(P\subset\R^d\), \(d\ge2\), be a convex body. Then \(\Pi P\) exists,
and for every invertible linear map \(L\colon\R^d\to\R^d\),
\[
 \Pi(LP)=|\det L|L^{-t}\Pi P,\qquad R_d(LP)=R_d(P),
\]
where \(L^{-t}\) denotes the inverse transpose. Translations also
preserve \(R_d\).
\end{lemma}

\begin{proof}
\emph{Plan.} Translations are immediate (Step~1). For a polytope we compute
how \(L\) changes the unit normal and the area of each facet (Steps~2--3).
Each area-normal vector is multiplied by \(M:=|\det L|L^{-t}\), so
\eqref{eq:zonotope} gives the claim (Step~4). A general body is squeezed
between a polytope and a slightly larger dilate of it (Step~5). This yields
the existence of \(\Pi P\) (Step~6) and lets the identity pass to the limit
(Steps~7--8).

\step{Step 1 (translations).} For \(t\in\R^d\),
\(\proj_{u^\perp}(P+t)\) is a translate of \(\proj_{u^\perp}P\). So \(P+t\)
and \(P\) have the same projection volumes, hence the same projection body.
Also \(|P+t|=|P|\), so \(R_d(P+t)=R_d(P)\).

\medskip\noindent\emph{Polytope case.} In Steps~2--4, \(P\) is a
full-dimensional polytope. Then \(LP\) is a polytope whose facets are the
images \(LF\) of the facets \(F\) of \(P\).

\step{Step 2 (normals).}
\emph{The outward unit normal of \(LF\) is
\(\nu_{LF}=L^{-t}\nu_F/\|L^{-t}\nu_F\|\).}
For \(y=Lx\),
\(\langle\nu_F,x\rangle=\langle\nu_F,L^{-1}y\rangle
=\langle L^{-t}\nu_F,y\rangle\). So the relations
\(\langle\nu_F,x\rangle=h_P(\nu_F)\) on \(F\) and
\(\langle\nu_F,x\rangle\le h_P(\nu_F)\) on \(P\) become
\(\langle L^{-t}\nu_F,y\rangle=h_P(\nu_F)\) on \(LF\) and \(\le\) on
\(LP\). Thus \(L^{-t}\nu_F\) is an outward normal of \(LF\), and we
normalize it. The argument does not involve \(\det L\), so it holds for
either sign of \(\det L\).

\step{Step 3 (areas).}
\emph{\(s_{LF}=|\det L|\,\|L^{-t}\nu_F\|\,s_F\).}
Let \(Y=F+[0,\nu_F]\) be the prism of height one over \(F\), so
\(|Y|=s_F\). Then \(LY=LF+[0,L\nu_F]\) and \(|LY|=|\det L|s_F\). Measure
height above the hyperplane of \(LF\) by
\(g(y)=\langle\nu_{LF},y\rangle-c\), where \(c\) is the constant value of
\(\langle\nu_{LF},\cdot\rangle\) on \(LF\). The set \(LY\) is the union of
the translates \(LF+tL\nu_F\), \(t\in[0,1]\), and \(g\equiv t\eta\) on the
translate with parameter \(t\), where
\[
 \eta=\langle L\nu_F,\nu_{LF}\rangle
 =\frac{\langle L\nu_F,L^{-t}\nu_F\rangle}{\|L^{-t}\nu_F\|}
 =\frac{\langle\nu_F,L^{t}L^{-t}\nu_F\rangle}{\|L^{-t}\nu_F\|}
 =\frac{1}{\|L^{-t}\nu_F\|}>0.
\]
Since \(\eta\ne0\), different \(t\) give different heights. So for each
\(h\in[0,\eta]\) the slice \(\{y\in LY:g(y)=h\}\) is exactly the translate
\(LF+(h/\eta)L\nu_F\), of \((d-1)\)-volume \(s_{LF}\), and the slice is
empty for other \(h\). Fubini in the direction \(\nu_{LF}\) gives
\(|LY|=\int_0^\eta s_{LF}\,dh=s_{LF}\,\eta\). Thus
\(s_{LF}=|\det L|s_F/\eta\), which is the claim.

\step{Step 4 (conclusion for polytopes).} By Steps~2 and~3,
\(s_{LF}\nu_{LF}=|\det L|L^{-t}(s_F\nu_F)=M(s_F\nu_F)\), also when
\(\det L<0\). A linear map sends segments to segments and commutes with
Minkowski sums, so \eqref{eq:zonotope} for \(LP\) and for \(P\) gives
\[
 \Pi(LP)=\sum_F\Bigl[-\tfrac12Ms_F\nu_F,\tfrac12Ms_F\nu_F\Bigr]
 =M\sum_F\Bigl[-\tfrac12s_F\nu_F,\tfrac12s_F\nu_F\Bigr]=M\,\Pi P.
\]
Now \(|\det M|=|\det L|^d\,|\det L|^{-1}=|\det L|^{d-1}\), so
\(|\Pi(LP)|=|\det L|^{d-1}|\Pi P|\). Also
\(|LP|^{d-1}=|\det L|^{d-1}|P|^{d-1}\). Dividing, \(R_d(LP)=R_d(P)\).

\medskip\noindent\emph{General case.} Now \(P\) is an arbitrary convex
body. By Step~1 we may translate \(P\) so that \(0\) is an interior point.
This changes neither side of the covariance identity, since
\(L(P+t)=LP+Lt\).

\step{Step 5 (approximating polytopes).} \emph{There are \(a>0\),
polytopes \(P_1\subset P_2\subset\cdots\) and numbers \(\varepsilon_m\to0\)
with}
\begin{equation}\label{eq:sandwich}
 aB_2^d\subset P_m\subset P\subset P_m+\varepsilon_mB_2^d\subset\mu_mP_m,
 \qquad\mu_m:=1+\varepsilon_m/a\longrightarrow1.
\end{equation}
\emph{Construction.} Choose \(a>0\) with \(P_0:=[-a,a]^d\subset P\); then
\(aB_2^d\subset P_0\). For each \(m\ge1\), compactness gives a finite set
\(N_m\subset P\) such that every point of \(P\) is within \(1/m\) of
\(N_m\). Put \(P_m=\conv(P_0\cup N_1\cup\cdots\cup N_m)\) and
\(\varepsilon_m=1/m\). Then \(P_m\) is the convex hull of finitely many
points, the sequence is increasing, and \(aB_2^d\subset P_m\subset P\) by
convexity of \(P\). Also \(P\subset N_m+\varepsilon_mB_2^d\subset
P_m+\varepsilon_mB_2^d\). Finally
\(\varepsilon_mB_2^d=(\varepsilon_m/a)\,aB_2^d\subset(\varepsilon_m/a)P_m\),
and \(P_m+(\varepsilon_m/a)P_m=\mu_mP_m\) because \(P_m\) is convex.

Fix also \(\rho>0\) with \(P\subset\rho B_2^d\).

\step{Step 6 (existence of \(\Pi P\)).} For unit \(u\) put
\(\varphi(u)=\vol_{d-1}(\proj_{u^\perp}P)\) and
\(\varphi_m(u)=\vol_{d-1}(\proj_{u^\perp}P_m)=h_{\Pi P_m}(u)\); here
\(\Pi P_m\) exists by Lemma~\ref{lem:facets}. Projecting
\eqref{eq:sandwich} gives
\(\proj_{u^\perp}P_m\subset\proj_{u^\perp}P\subset
\mu_m\proj_{u^\perp}P_m\). Taking \((d-1)\)-volumes,
\begin{equation}\label{eq:phi}
 \varphi_m\le\varphi\le\mu_m^{d-1}\varphi_m,
 \qquad\text{so}\qquad
 0\le\varphi-\varphi_m
 \le(\mu_m^{d-1}-1)\kappa_{d-1}\rho^{d-1}.
\end{equation}
The last bound uses
\(\varphi_m(u)\le\vol_{d-1}(\proj_{u^\perp}\rho B_2^d)
=\kappa_{d-1}\rho^{d-1}\).
Thus \(\varphi_m\to\varphi\) uniformly on \(S^{d-1}\). Moreover:
\begin{itemize}[leftmargin=2em,itemsep=2pt]
\item \(\varphi_m\le\varphi_{m+1}\) because \(P_m\subset P_{m+1}\), so
 \(\Pi P_m\subset\Pi P_{m+1}\) by (S4).
\item \(\kappa_{d-1}a^{d-1}\le\varphi_m\le\kappa_{d-1}\rho^{d-1}\) on
 \(S^{d-1}\) because \(aB_2^d\subset P_m\subset\rho B_2^d\). By (S3) and
 (S4),
 \(\kappa_{d-1}a^{d-1}B_2^d\subset\Pi P_m
 \subset\kappa_{d-1}\rho^{d-1}B_2^d\).
\end{itemize}
Let \(C\) be the closure of \(\bigcup_m\Pi P_m\). An increasing union of
convex sets is convex, and so is its closure. With the second bullet, \(C\)
is a compact convex set with nonempty interior. The supremum of
\(\langle u,\cdot\rangle\) over a set equals its supremum over the closure,
and the \(\varphi_m\) increase. So for unit \(u\),
\(h_C(u)=\sup_mh_{\Pi P_m}(u)=\lim_m\varphi_m(u)=\varphi(u)\). Thus \(C\)
has the support function required of \(\Pi P\), that is, \(\Pi P\) exists
and equals \(C\). By \eqref{eq:phi} and (S4),
\begin{equation}\label{eq:volsandwich}
 \Pi P_m\subset\Pi P\subset\mu_m^{d-1}\Pi P_m,
 \qquad
 |\Pi P_m|\le|\Pi P|\le\mu_m^{d(d-1)}|\Pi P_m|,
 \qquad
 |P_m|\le|P|\le\mu_m^d|P_m|.
\end{equation}
The last chain comes from \eqref{eq:sandwich}.

\step{Step 7 (covariance).} Apply \(L\) to \eqref{eq:sandwich}:
\(LP_m\subset LP\subset\mu_mLP_m\), with polytopes \(LP_m\). Let
\(\psi(u)=\vol_{d-1}(\proj_{u^\perp}LP)\) for unit \(u\). Exactly as for
\eqref{eq:phi},
\(h_{\Pi(LP_m)}(u)\le\psi(u)\le\mu_m^{d-1}h_{\Pi(LP_m)}(u)\), so
\(h_{\Pi(LP_m)}(u)\to\psi(u)\). On the other hand, Step~4 and (S2) give
\[
 h_{\Pi(LP_m)}(u)=h_{M\Pi P_m}(u)=h_{\Pi P_m}(M^tu)
 \longrightarrow h_{\Pi P}(M^tu)=h_{M\Pi P}(u).
\]
The convergence holds because \(h_{\Pi P_m}\to h_{\Pi P}\) on \(S^{d-1}\)
by Step~6, hence on \(\R^d\) by homogeneity. Therefore
\(\psi=h_{M\Pi P}\) on \(S^{d-1}\). So \(M\Pi P\) is the projection body
of \(LP\): \(\Pi(LP)=|\det L|L^{-t}\Pi P\).

\step{Step 8 (invariance of \(R_d\)).} Two arguments are available; each is
complete on its own.

\emph{(a) The source's argument.} By \eqref{eq:volsandwich},
\(\mu_m^{-d(d-1)}R_d(P_m)\le R_d(P)\le\mu_m^{d(d-1)}R_d(P_m)\), so
\(R_d(P_m)\to R_d(P)\). The sandwich \(LP_m\subset LP\subset\mu_mLP_m\) and
the inequalities for \(\psi\) in Step~7 give \eqref{eq:volsandwich} with
\(LP_m,LP\) in place of \(P_m,P\) in the same way, so
\(R_d(LP_m)\to R_d(LP)\). Since \(R_d(LP_m)=R_d(P_m)\) by Step~4,
\(R_d(LP)=R_d(P)\).

\emph{(b) A direct cross-check.} Step~7 gives \(\Pi(LP)=M\,\Pi P\) for the
convex body \(P\). Hence
\(|\Pi(LP)|=|\det M|\,|\Pi P|=|\det L|^{d-1}|\Pi P|\), and
\(|LP|^{d-1}=|\det L|^{d-1}|P|^{d-1}\). These are the only facts used at
the end of Step~4, and none of them requires \(P\) to be a polytope.
\end{proof}

\section{The simplex value}\label{sec:simplex}

We now compute the value to which the product will be compared. The
projection body of a simplex is a coordinate cube plus one extra segment.
The proposition is used in the proofs of Theorem~\ref{thm:counterexample}
and Corollary~\ref{cor:exponential-excess}, which need only the first
sentence, for the standard simplex.

\begin{proposition}\label{prop:simplex}
For \(d\ge2\), the standard simplex
\(T_d=\conv(0,e_1,\ldots,e_d)\subset\R^d\) satisfies
\[
 |T_d|=\frac1{d!},\qquad
 |\Pi T_d|=\frac{d+1}{((d-1)!)^d},\qquad
 R_d(T_d)=c_d.
\]
Every full-dimensional \(d\)-simplex has the same value \(c_d\).
\end{proposition}

\begin{proof}
\emph{Plan.} Compute the \(d+1\) area-normal vectors of \(T_d\)
(Steps~1--2). Recognize \(\Pi T_d\) via \eqref{eq:zonotope} as a translate
of \(([0,1]^d+[0,w])/(d-1)!\) with \(w=e_1+\dots+e_d\) (Step~3). Compute
the volume of \([0,1]^d+[0,w]\) by slicing along \(w\) (Steps~4--5), and
assemble (Steps~6--7).

\step{Step 1 (volume).} \emph{\(|T_j|=1/j!\) for all \(j\ge1\).} This holds
for \(T_1=[0,1]\). For \(j\ge2\), the slice of \(T_j\) at height
\(x_j=t\in[0,1]\) is a copy of \((1-t)T_{j-1}\), so
\(|T_j|=\int_0^1(1-t)^{j-1}|T_{j-1}|\,dt=|T_{j-1}|/j\). Induction gives the
claim.

\step{Step 2 (area-normal vectors).} We have
\(T_d=\{x:x_i\ge0\ (1\le i\le d),\ x_1+\dots+x_d\le1\}\), with \(d+1\)
facets.
\begin{itemize}[leftmargin=2em,itemsep=2pt]
\item \emph{Coordinate facets.} The facet in \(\{x_i=0\}\) is
 \(\conv(0,e_j:j\ne i)\), a copy of \(T_{d-1}\). Its outward unit normal is
 \(-e_i\) and its area is \(1/(d-1)!\) by Step~1. Area-normal vector:
 \(-e_i/(d-1)!\).
\item \emph{Slanted facet.} \(F_0=\conv(e_1,\dots,e_d)\) lies in
 \(\{x_1+\dots+x_d=1\}\) and has outward unit normal \(w/\sqrt d\).
 Dropping the last coordinate is the orthogonal projection onto
 \(e_d^\perp\). It maps \(F_0\) bijectively onto the standard simplex
 \(T_{d-1}\subset\R^{d-1}\), with inverse
 \(x'\mapsto(x',1-\sum_ix'_i)\). By Step~1 of Lemma~\ref{lem:facets}, with
 \(u=e_d\), this projection multiplies areas on the hyperplane of \(F_0\)
 by \(|\langle w/\sqrt d,e_d\rangle|=1/\sqrt d\). Hence
 \(1/(d-1)!=s_{F_0}/\sqrt d\), that is, \(s_{F_0}=\sqrt d/(d-1)!\).
 Area-normal vector:
 \(\dfrac{\sqrt d}{(d-1)!}\cdot\dfrac{w}{\sqrt d}=\dfrac{w}{(d-1)!}\).
\end{itemize}
Sanity check: the \(d+1\) area-normal vectors sum to
\((-e_1-\dots-e_d+w)/(d-1)!=0\), as they must for a closed polytope.

\step{Step 3 (\(\Pi T_d\) as a zonotope).} Put \(\theta=1/(d-1)!\). By
\eqref{eq:zonotope},
\[
 \Pi T_d=\sum_{i=1}^d\Bigl[-\tfrac\theta2e_i,\tfrac\theta2e_i\Bigr]
 +\Bigl[-\tfrac\theta2w,\tfrac\theta2w\Bigr].
\]
The sign of the generator \(-e_i\) is irrelevant, because a centered
segment is unchanged when its generator is negated. Since
\([-\tfrac12v,\tfrac12v]=-\tfrac12v+[0,v]\) and
\(\sum_i[0,e_i]=[0,1]^d\), the body \(\Pi T_d\) is a translate of
\[
 \frac{[0,1]^d+[0,w]}{(d-1)!}.
\]

\step{Step 4 (extrusion volume).} Let \(Q=[0,1]^d\). \emph{Then
\(|Q+[0,w]|=1+h_{\Pi Q}(w)\).} Let \(\hat w=w/\|w\|\) and
\(\pi=\proj_{w^\perp}\). For a set \(S\) and \(y\in w^\perp\), let
\(S_y=\{t\in\R:y+t\hat w\in S\}\) be its fiber over \(y\). Note that
\([0,w]=\{\sigma\hat w:0\le\sigma\le\|w\|\}\).
\begin{itemize}[leftmargin=2em,itemsep=2pt]
\item If \(y\in\pi(Q)\), then \(Q_y\) is a nonempty compact interval
 \([t_0,t_1]\), possibly a single point, because \(Q\) is compact and
 convex. A point \(y+t\hat w\) lies in \(Q+[0,w]\) if and only if
 \(y+t\hat w=x+\sigma\hat w\) with \(x\in Q\) and
 \(\sigma\in[0,\|w\|]\), that is, if and only if
 \(t-\sigma\in[t_0,t_1]\) for some \(\sigma\in[0,\|w\|]\). So
 \((Q+[0,w])_y=[t_0,t_1+\|w\|]\), of length \((t_1-t_0)+\|w\|\).
\item If \(y\notin\pi(Q)\), then \((Q+[0,w])_y=\emptyset\), because
 \(\pi(w)=0\) gives \(\pi(Q+[0,w])=\pi(Q)\).
\end{itemize}
By Fubini, applied to \(Q+[0,w]\) and to \(Q\),
\[
 |Q+[0,w]|
 =\int_{\pi(Q)}\bigl(\operatorname{length}(Q_y)+\|w\|\bigr)\,dy
 =|Q|+\|w\|\vol_{d-1}(\proj_{w^\perp}Q)
 =1+h_{\Pi Q}(w).
\]
The last equality uses \(|Q|=1\), the definition of \(\Pi Q\) at the unit
vector \(\hat w\), and homogeneity: \(\|w\|h_{\Pi Q}(\hat w)=h_{\Pi Q}(w)\).

\step{Step 5 (\(h_{\Pi Q}(w)=d\)).} The cube has \(2d\) facets, with unit
normals \(\pm e_i\) and area one, and \(|\langle\pm e_i,w\rangle|=1\). By
\eqref{eq:facet}, \(h_{\Pi Q}(w)=\tfrac12\cdot2d\cdot1=d\). Hence
\(|Q+[0,w]|=d+1\). Figure~\ref{fig:extrusion} shows the case \(d=2\).

\step{Step 6 (assembling).} Volume is translation invariant, and dilation
by \(1/(d-1)!\) in \(\R^d\) multiplies volume by \(((d-1)!)^{-d}\). So
Steps~3--5 give \(|\Pi T_d|=(d+1)/((d-1)!)^d\). With \(|T_d|=1/d!\) and
\((d-1)!=d!/d\),
\[
 R_d(T_d)=\frac{(d+1)(d!)^{d-1}}{((d-1)!)^d}
 =\frac{(d+1)(d!)^{d-1}d^d}{(d!)^d}
 =\frac{(d+1)d^d}{d!}=c_d.
\]
Sanity check for \(d=2\): \(T_2\) has area \(\tfrac12\), \(\Pi T_2\) is a
translate of the hexagon of Figure~\ref{fig:extrusion} with area \(3\), and
\(3/\tfrac12=6=c_2\).

\step{Step 7 (all simplices).} A full-dimensional simplex
\(\conv(v_0,\dots,v_d)\) equals \(v_0+LT_d\), where \(L\) is the linear map
with \(Le_i=v_i-v_0\). This map is invertible because the \(v_i\) are
affinely independent. By Lemma~\ref{lem:affine} (polytope case, Steps~1
and~4), \(R_d(v_0+LT_d)=R_d(T_d)=c_d\).
\end{proof}

\begin{figure}[tb]
\centering
\begin{tikzpicture}[scale=1.5, line join=round]
 \fill[black!7] (0,0) rectangle (1,1);
 \fill[blue!13] (1,0)--(2,1)--(2,2)--(1,1)--cycle;
 \fill[yellow!23] (0,1)--(1,1)--(2,2)--(1,2)--cycle;
 \draw[thick] (0,0)--(1,0)--(2,1)--(2,2)--(1,2)--(0,1)--cycle;
 \draw[thin] (0,1)--(1,1)--(1,0);
 \draw[thin] (1,1)--(2,2);
 \node at (.5,.5) {\(Q\)};
 \node at (1.5,1) {\(1\)};
 \node at (1,1.5) {\(1\)};
 \node[below left] at (0,0) {\(0\)};
 \node[above right] at (2,2) {\((2,2)\)};
 \draw[->,thick] (2.6,.3)--(3.6,1.3);
 \node[right] at (3.15,.65) {\(w=(1,1)\)};
\end{tikzpicture}
\caption{The planar extrusion \(Q+[0,(1,1)]\), with
\(Q=[0,1]^2\), consists of the unit square and two parallelograms of
area one. Its area is \(3\). The fiber argument in
Proposition~\ref{prop:simplex} proves the corresponding volume \(d+1\)
in every dimension.}
\label{fig:extrusion}
\end{figure}

\section{The counterexample and exponential excess}\label{sec:counterexample}

The geometric work is complete. What remains is arithmetic with the
constants \(c_d\).

\begin{proof}[Proof of Theorem~\ref{thm:counterexample}]
\step{Step 1.} \(T_{10}\) is a full-dimensional convex polytope in
\(\R^{10}\), and \(10\ge2\). So \(K=T_{10}\times T_{10}\) is a convex body
in \(\R^{20}\), and Proposition~\ref{prop:product} with \(A=B=T_{10}\),
followed by Proposition~\ref{prop:simplex}, gives
\[
 R_{20}(K)=R_{10}(T_{10})^2=c_{10}^2
         =\left(\frac{11\cdot10^{10}}{10!}\right)^2.
\]

\step{Step 2.} Dividing by \(c_{20}=21\cdot20^{20}/20!\) and using
\(10^{20}/20^{20}=2^{-20}\) and \(20!/(10!)^2=\binom{20}{10}\),
\[
 \frac{R_{20}(K)}{c_{20}}
 =\frac{121\cdot10^{20}\cdot20!}
        {(10!)^2\cdot21\cdot20^{20}}
 =\frac{121\binom{20}{10}}{21\cdot2^{20}}.
\]

\step{Step 3.} \(\binom{20}{10}=184\,756\) and \(2^{20}=1\,048\,576\).
Then \(121\cdot184\,756=22\,355\,476\) and
\(21\cdot1\,048\,576=22\,020\,096\), so
\[
 121\cdot184\,756-21\cdot1\,048\,576=335\,380>0.
\]
The ratio is therefore \(22\,355\,476/22\,020\,096>1\) (about \(1.015\)).
That is, \(|\Pi K|>c_{20}|K|^{19}\), a strict violation of
\eqref{eq:proposed-bound}.
\end{proof}

\noindent\textbf{The general product test.} For \(r,s\ge2\), the same two
propositions give \(R_{r+s}(T_r\times T_s)=c_rc_s\). Dividing by
\(c_{r+s}\) and using \((r+s)!/(r!\,s!)=\binom{r+s}{r}\),
\[
 \frac{R_{r+s}(T_r\times T_s)}{c_{r+s}}
 =\frac{c_rc_s}{c_{r+s}}
 =\frac{(r+1)(s+1)}{r+s+1}
   \binom{r+s}{r}\frac{r^rs^s}{(r+s)^{r+s}}.
\]
Since \(T_r\times T_s\) is a convex body in \(\R^{r+s}\), any universal
upper constant in dimension \(r+s\) must be at least \(c_rc_s\). The case
\(r=s=10\) is Theorem~\ref{thm:counterexample}.

Iterating the product identity gives the following quantitative
consequence.

\begin{corollary}\label{cor:exponential-excess}
There exist a constant \(\lambda>1\) and an integer \(N\ge20\) such that,
for every integer \(n\ge N\), some full-dimensional convex polytope
\(K_n\subset\R^n\) satisfies
\[
 \frac{R_n(K_n)}{R_n(T_n)}\ge\lambda^n.
\]
The polytopes \(K_n\) may be chosen as Cartesian products of simplices.
\end{corollary}

\begin{proof}
\emph{Plan.} Take \(k\) factors \(T_{10}\) and one more simplex of
dimension between \(10\) and \(19\) to reach dimension \(n\) (Step~1).
Bound \(c_n\) above by \((n+1)\beta^n\) with \(\beta=11/4\) (Step~2). Check
that \(c_{10}>\beta^{10}\), so that each ten-dimensional factor
contributes more than \(\beta^{10}\) (Step~3). Then compare (Steps~4--5).

\step{Step 1 (the bodies).} For \(r\in\{0,\ldots,9\}\) put \(d_r=10+r\).
Every integer \(n\ge20\) has a unique representation \(n=10k+d_r\) with
\(k\ge1\) and \(r\in\{0,\dots,9\}\): divide \(n-10\ge10\) by \(10\) with
remainder \(r\). Let
\[
 K_n=T_{10}^{\,k}\times T_{d_r}\subset\R^n,
\]
where \(T_{10}^{\,k}\) is the \(k\)-fold Cartesian product. A product of
full-dimensional convex polytopes is one, and all factor dimensions are at
least \(10\ge2\). So Proposition~\ref{prop:product} applies \(k\) times,
and with Proposition~\ref{prop:simplex},
\[
 R_n(K_n)=c_{10}^k\,c_{d_r},\qquad R_n(T_n)=c_n.
\]

\step{Step 2 (upper bound for \(c_n\)).} Put \(\beta=11/4\). \emph{Then
\(e<\beta\) and \(c_n<(n+1)\beta^n\) for \(n\ge2\).} For \(i\ge0\) we have
\((3+i)!=6\cdot\prod_{m=4}^{3+i}m\ge6\cdot4^i\): the product has \(i\)
factors, each at least \(4\), and is empty for \(i=0\). Hence
\[
 \frac1{(3+i)!}\le\frac16\,4^{-i}\le\frac16\,3^{-i},
\]
and the second inequality is strict for \(i\ge1\). (This is the source's
remark that successive ratios in the tail are less than \(1/3\).) Summing
over \(i\ge0\),
\[
 e=\frac52+\sum_{j=3}^{\infty}\frac1{j!}
 <\frac52+\frac16\sum_{i=0}^{\infty}3^{-i}
 =\frac52+\frac16\cdot\frac32=\frac{11}4=\beta.
\]
Next, \(n^n/n!\) is one term of the series
\(e^n=\sum_{j\ge0}n^j/j!\), whose terms are all positive. So
\(n^n/n!<e^n\) and
\[
 c_n=(n+1)\frac{n^n}{n!}<(n+1)e^n<(n+1)\beta^n.
\]

\step{Step 3 (the ten-dimensional block).} We have \(10!=3\,628\,800\), so
\(10^{10}/10!=10^8/36\,288\). With \(36\,288=2^6\cdot567\) and
\(10^8=2^8\cdot5^8\), this is \(4\cdot5^8/567=1\,562\,500/567\). Hence
\[
 c_{10}=\frac{11\cdot10^{10}}{10!}=\frac{17\,187\,500}{567},
 \qquad
 \alpha:=\frac{c_{10}}{\beta^{10}}
 =\frac{2^2\cdot5^8\cdot11\cdot4^{10}}{567\cdot11^{10}}
 =\frac{2^{22}\cdot5^8}{567\cdot11^9}
 =\frac{1\,638\,400\,000\,000}{1\,336\,956\,340\,797}>1.
\]
Here \(2^{22}\cdot5^8=2^{14}\cdot10^8\) and \(2^{14}=16\,384\). Also
\(11^9=2\,357\,947\,691\) and
\(567\cdot2\,357\,947\,691=1\,336\,956\,340\,797\). So
\(\alpha\approx1.2255\).

\step{Step 4 (comparison).} Fix \(\lambda\) with
\(1<\lambda<\alpha^{1/10}\), so that \(\alpha/\lambda^{10}>1\). For
\(n=10k+d_r\), Steps~1 and~2 give
\[
 \frac{R_n(K_n)}{R_n(T_n)\lambda^n}
 =\frac{c_{10}^k\,c_{d_r}}{c_n\lambda^n}
 >\frac{c_{10}^k\,c_{d_r}}{(n+1)\beta^n\lambda^n}.
\]
Substitute \(c_{10}^k=\alpha^k\beta^{10k}\),
\(\beta^n=\beta^{10k}\beta^{d_r}\),
\(\lambda^n=\lambda^{10k}\lambda^{d_r}\) and \(n+1=10k+d_r+1\). The factor
\(\beta^{10k}\) cancels, and
\[
 \frac{R_n(K_n)}{R_n(T_n)\lambda^n}
 >c_{d_r}\beta^{-d_r}\lambda^{-d_r}\,
   \frac{(\alpha/\lambda^{10})^k}{10k+d_r+1}.
\]

\step{Step 5 (thresholds).} Fix \(r\). The right side is a positive
constant times \(\gamma^k/(10k+d_r+1)\) with
\(\gamma=\alpha/\lambda^{10}>1\). It tends to infinity as \(k\to\infty\),
since an exponential with base \(>1\) outgrows a linear function. So there
is \(k_r\ge1\) such that the right side is at least \(1\) for all
\(k\ge k_r\). Let \(N=\max_{0\le r\le9}(10k_r+d_r)\); then \(N\ge20\). If
\(n\ge N\) and \(n=10k+d_r\), then \(10k+d_r\ge10k_r+d_r\), so
\(k\ge k_r\). Step~4 then gives \(R_n(K_n)/R_n(T_n)>\lambda^n\).
\end{proof}

\section*{Verification checklist}

No step of this rewrite is marked as unverified. A sceptical reader should
check the following load-bearing points first.
\begin{enumerate}[leftmargin=2em,itemsep=3pt]
\item \textbf{The factor \(\tfrac12\) in the facet formula.} Almost every
 line parallel to \(u\) through \(P\) meets exactly two facets, so
 \(\sum_Fs_F|\langle\nu_F,u\rangle|=2\vol_{d-1}(\proj_{u^\perp}P)\).
 See Lemma~\ref{lem:facets}, Steps~1--3.
\item \textbf{Facets of a product and the weights in
 \eqref{eq:product-body}.} The facets of \(A\times B\) are exactly
 \(F\times B\) and \(A\times G\), with areas \(s_F|B|\) and \(|A|s_G\).
 The exponent bookkeeping
 \(|B|^r|A|^s/(|A||B|)^{r+s-1}=|A|^{-(r-1)}|B|^{-(s-1)}\) must come out.
 See Proposition~\ref{prop:product}, Steps~1, 2 and~4.
\item \textbf{The slanted facet of the simplex.} Its area is
 \(\sqrt d/(d-1)!\), so its area-normal vector is \(w/(d-1)!\), and the
 \(d+1\) area-normal vectors sum to zero. See
 Proposition~\ref{prop:simplex}, Step~2.
\item \textbf{The extrusion volume \(|[0,1]^d+[0,w]|=d+1\).} It rests on
 the fiber-length identity and \(h_{\Pi Q}(w)=d\). See
 Proposition~\ref{prop:simplex}, Steps~4--5; test case \(d=2\) in
 Figure~\ref{fig:extrusion} and Step~6.
\item \textbf{The arithmetic of Theorem~\ref{thm:counterexample}.}
 \(c_{10}^2/c_{20}=121\binom{20}{10}/(21\cdot2^{20})\) and
 \(121\cdot184\,756-21\cdot1\,048\,576=335\,380>0\). See the proof of
 Theorem~\ref{thm:counterexample}, Steps~2--3.
\item \textbf{Dependence.} Theorem~\ref{thm:counterexample} needs only
 items 1--5. Lemma~\ref{lem:affine} enters only through ``every simplex
 has value \(c_d\)'' (Proposition~\ref{prop:simplex}, Step~7), and there
 only in its polytope case. Its key computation is the prism height
 \(\langle L\nu_F,\nu_{LF}\rangle=1/\|L^{-t}\nu_F\|\) in Step~3 of
 Lemma~\ref{lem:affine}.
\item \textbf{For Corollary~\ref{cor:exponential-excess}.} Check
 \(e<11/4\) and \(c_n<(n+1)(11/4)^n\) (Step~2), the constant
 \(\alpha=2^{22}5^8/(567\cdot11^9)>1\) (Step~3), and the cancellation of
 \(\beta^{10k}\) (Step~4).
\end{enumerate}

\begin{thebibliography}{7}

\bibitem{Brannen}
N.~S.~Brannen,
\emph{Volumes of projection bodies},
Mathematika \textbf{43} (1996), no.~2, 255--264.
\href{https://doi.org/10.1112/S002557930001175X}{doi:10.1112/S002557930001175X}.

\bibitem{ChenEtAl}
S.~Chen, Y.~Feng, Y.~Li, D.~Xi, and L.~Xu,
\emph{Petty's conjectured projection inequality in dimension three},
preprint, 2026, MathSciDoc:2608.23001.
\href{https://archive.ymsc.tsinghua.edu.cn/pacm_download/741/12780-pettyinequality.pdf}{Full text}.

\bibitem{FengEtAl}
Y.~Feng, S.~Hu, W.~Liu, and L.~Xu,
\emph{On the reverse projection inequality},
preprint, 2026, MathSciDoc:2608.23002.
\href{https://archive.ymsc.tsinghua.edu.cn/pacm_download/743/12781--2026.8.26.pdf}{Full text}.

\bibitem{Henk}
M.~Henk,
\emph{Note on projection bodies of zonotopes with \(n+1\) generators},
Acta Mathematica Scientia \textbf{45} (2025), 96--103.
\href{https://doi.org/10.1007/s10473-025-0107-9}{doi:10.1007/s10473-025-0107-9}.

\bibitem{Ludwig}
M.~Ludwig,
\emph{Projection bodies and valuations},
Advances in Mathematics \textbf{172} (2002), no.~2, 158--168.
\href{https://doi.org/10.1016/S0001-8708(02)00021-X}{doi:10.1016/S0001-8708(02)00021-X}.

\bibitem{Petty}
C.~M.~Petty,
\emph{Projection bodies},
in \emph{Proceedings of the Colloquium on Convexity (Copenhagen, 1965)},
K\o benhavns Universitet Matematiske Institut, Copenhagen, 1967,
pp.~234--241.

\bibitem{Saroglou}
C.~Saroglou,
\emph{On the shape of a convex body with respect to its second projection body},
Advances in Applied Mathematics \textbf{67} (2015), 55--74.
\href{https://doi.org/10.1016/j.aam.2015.03.004}{doi:10.1016/j.aam.2015.03.004}.
Preprint:
\href{https://arxiv.org/abs/1409.4347v2}{arXiv:1409.4347v2}.

\end{thebibliography}

\end{document}
